% Tentative ML model proposal - one-page, two-column format.
% Compile with XeLaTeX or pdfLaTeX.
\documentclass[11pt,letterpaper]{article}
\usepackage[letterpaper,top=0.68in,bottom=0.68in,left=0.7in,right=0.7in,columnsep=0.23in]{geometry}
\usepackage{iftex}
\ifPDFTeX
  \usepackage[T1]{fontenc}
  \usepackage[utf8]{inputenc}
  \usepackage{mathptmx}
\else
  \usepackage{fontspec}
  \IfFontExistsTF{Times New Roman}{\setmainfont{Times New Roman}}{\IfFontExistsTF{Tinos}{\setmainfont{Tinos}}{\setmainfont{Liberation Serif}}}
\fi
\usepackage{microtype}
\usepackage{titlesec}
\usepackage{multicol}
\usepackage{hyperref}
\hypersetup{hidelinks}
\pagestyle{empty}
\setlength{\parindent}{0pt}
\setlength{\parskip}{3pt}
\titleformat{\section}{\normalfont\normalsize}{\Roman{section}.}{0.5em}{}
\titlespacing*{\section}{0pt}{7pt}{2pt}

\begin{document}
\begin{center}
{\fontsize{14}{16}\selectfont Tentative Machine Learning Model for Random Number Generator Data}\par
\vspace{3pt}
{\normalsize CPE 486/586 -- Machine Learning for Engineers}\par
\vspace{8pt}
\end{center}
\begin{multicols}{2}

\section{Problem and Objective}
Random number generators (RNGs) are used in simulation, security, and computing systems. Their outputs should have good statistical properties and should be difficult to predict. For this project, I want to investigate whether machine learning can find patterns in binary sequences generated by different sources. The main task will be to identify the source group of a sequence. I also plan to investigate whether previous bits can help predict the next bit.

\section{Dataset and Features}
I will use a public dataset from the Noise versus Randomness project. It contains 6,000 sequences of 100 bits each, with source labels 1, 2, and 3. There are 2,000 sequences for each label. Since I have not verified which physical device or simulator belongs to each label, I will treat them as three source groups.

For the source classification task, I have started with three features: frequency of ones, transition rate, and longest run of equal bits. I plan to expand the feature set to around 13 statistical measurements, including bit-pair frequencies, lag correlation, run statistics, and differences between smaller blocks of a sequence. For the next-bit task, I will use the previous 16 bits as input and the following bit (0 or 1) as the target.

\section{Proposed Models}
I will first use a simple baseline that predicts the most common label. Then I will compare logistic regression and random forest. Logistic regression will provide a basic linear model, while random forest can learn more complicated relationships between features. If time permits, I will also test an RBF support vector machine (SVM).

I will use standardization for models that need it and try principal component analysis (PCA) to reduce the number of input features. I will compare classification results with and without PCA. For next-bit prediction, I will begin with a majority-bit baseline and logistic regression, and then try random forest if possible.

\section{Training and Evaluation}
I plan to split the original sequences into 80\% training data and 20\% testing data, keeping the three source groups balanced. I will use five-fold cross-validation on the training set when comparing models. Scaling and PCA will be fitted using training data only. For next-bit prediction, I will split the original sequences before creating sliding windows so that overlapping windows from the same sequence do not appear in both sets.

For source classification, I will report accuracy, precision, recall, F1-score, and a confusion matrix. For next-bit prediction, I will use accuracy, F1-score, and log loss. I will also plot feature distributions and PCA projections to understand the data. I plan to calculate empirical Shannon entropy from bit frequencies and compare it with the prediction results. However, classification accuracy or a low prediction error alone cannot prove that an RNG is secure or truly random.

\section{Expected Outcome and Tools}
I expect this study to show whether simple statistical features are useful for separating the source groups and whether the previous bits contain information that helps predict the next bit. It is also possible that the models will perform close to the baseline. Since the dataset comes from another project, I will discuss the limits of the available source information.

I will use Python, scikit-learn, matplotlib, and VS Code. The dataset is available at \url{https://github.com/hubertkolcz/NoiseVsRandomness}. My code and results will be saved in my GitHub repository.

\end{multicols}
\end{document}
